\subsection{Interpolimi i Lagranzhit dhe termi i mbetjes}
Për nyje të ndryshme $x_0,\ldots,x_n$, polinomi interpolues është
\begin{equation}
 p_n(x)=\sum_{i=0}^{n}f(x_i)L_i(x),\qquad
 L_i(x)=\prod_{j\ne i}\frac{x-x_j}{x_i-x_j}.
\end{equation}
Vërtetimi është i drejtpërdrejtë: $L_i(x_j)=\delta_{ij}$, prandaj $p_n(x_j)=f(x_j)$. Uniciteti rrjedh sepse diferenca e dy polinomeve interpoluese do të kishte $n+1$ rrënjë, por gradë jo më të madhe se $n$.

Nëse $f\in C^{n+1}$, gabimi në një pikë $x$ është
\begin{equation}
 f(x)-p_n(x)=\frac{f^{(n+1)}(\xi_x)}{(n+1)!}
 \omega_{n+1}(x),\qquad
 \omega_{n+1}(x)=\prod_{i=0}^{n}(x-x_i).
\end{equation}
Kjo formulë ndan dy burime: lëmueshmërinë e funksionit dhe gjeometrinë e nyjeve. Nyjet e baraslarguara mund ta bëjnë $\omega_{n+1}$ shumë të madhe pranë skajeve; nyjet e Çebishevit e kontrollojnë këtë faktor dhe zbusin dukurinë Runge.

Forma baricentrike
\begin{equation}
 p_n(x)=\frac{\sum_i w_i f_i/(x-x_i)}{\sum_i w_i/(x-x_i)},
 \qquad w_i=\left[\prod_{j\ne i}(x_i-x_j)\right]^{-1},
\end{equation}
është numerikisht më e përshtatshme se zgjerimi i polinomit në fuqi. Nëse $x$ përkon me një nyje, kthehet drejtpërdrejt $f_i$ për të shmangur pjesëtimin me zero.

\subsection{Interpolimi Hermitian dhe ekstrapolimi i Richardson-it}
Interpolimi Hermitian përdor jo vetëm $f(x_i)$, por edhe derivate. Për dy nyje dhe vlera të derivateve ndërtohet një polinom kub. Në formë të përgjithshme, nyjet përsëriten në tabelën e diferencave të pjesëtuara dhe diferenca e parë në një nyje të përsëritur zëvendësohet nga $f'(x_i)$.

Ekstrapolimi nis nga një përafrim me zhvillim asimptotik
\begin{equation}
 A(h)=A+c_ph^p+c_{p+1}h^{p+1}+\cdots.
\end{equation}
Dy vlera me hapa $h$ dhe $h/r$ japin
\begin{equation}
 A_{\mathrm{R}}=\frac{r^pA(h/r)-A(h)}{r^p-1}
 =A+\Oh(h^{p+1}),
\end{equation}
ose rend edhe më të lartë kur zhvillimi përmban vetëm fuqi çifte. Kjo është ideja pas derivimit të përmirësuar dhe integrimit Romberg. Ekstrapolimi nuk duhet zbatuar jashtë regjimit asimptotik; nëse hapi është shumë i madh, modeli i gabimit nuk vlen, ndërsa për hap tepër të vogël dominon rrumbullakimi.

\subsection{Katrorët më të vegjël dhe QR}
Për modelin linear $\vect y=X\vect\beta+\vect\varepsilon$, minimizohet
\begin{equation}
 S(\vect\beta)=\|\vect y-X\vect\beta\|_2^2.
\end{equation}
Gradienti është $\nabla S=-2X^T(\vect y-X\vect\beta)$. Kushti $\nabla S=0$ jep ekuacionet normale
\begin{equation}
 X^TX\widehat{\vect\beta}=X^T\vect y.
\end{equation}
Gjeometrikisht, mbetja është ortogonale me hapësirën e kolonave të $X$. Megjithatë, $\kappa(X^TX)=\kappa(X)^2$, prandaj formimi i ekuacioneve normale përkeqëson kushtëzimin. Me $X=QR$, ku $Q^TQ=I$, zgjidhet
\begin{equation}
 R\widehat{\vect\beta}=Q^T\vect y,
\end{equation}
pa katrorëzuar numrin e kushtëzimit.

Nëse gabimet kanë kovariancë $\Sigma$, kriteri i peshuar është
\begin{equation}
 (\vect y-X\vect\beta)^T\Sigma^{-1}(\vect y-X\vect\beta).
\end{equation}
Kjo formulë tregon se peshat nuk janë dekorim: ato përfaqësojnë modelin probabilitar të matjeve. Kur gabimet kanë shpërndarje Gauss, katrorët më të vegjël përkojnë me vlerësimin e gjasës më të madhe.

\subsection{Shmangiet minimale absolute}
Kriteri
\begin{equation}
 \min_{\vect\beta}\sum_{i=1}^{n}|y_i-\vect x_i^T\vect\beta|
\end{equation}
është më pak i ndjeshëm ndaj vlerave ekstreme se kriteri kuadratik. Duke futur $u_i\ge0$ me
\begin{equation}
 -u_i\le y_i-\vect x_i^T\vect\beta\le u_i,
\end{equation}
problemi bëhet programim linear me objektiv $\min\sum_i u_i$. Funksioni absolut nuk është i derivueshëm në zero, prandaj nuk duhet trajtuar verbërisht me ekuacionet normale.

\begin{lstlisting}[language=Python,caption={Përshtatje lineare me QR dhe diagnostikë të mbetjeve.}]
import numpy as np

def pershtatje_qr(x, y):
    x = np.asarray(x, float); y = np.asarray(y, float)
    X = np.column_stack((np.ones_like(x), x))
    Q, R = np.linalg.qr(X, mode="reduced")
    beta = np.linalg.solve(R, Q.T @ y)
    mbetje = y - X @ beta
    sigma2 = (mbetje @ mbetje) / (len(y) - X.shape[1])
    kov_beta = sigma2 * np.linalg.inv(R) @ np.linalg.inv(R).T
    return beta, mbetje, kov_beta
\end{lstlisting}

\subsection{Përmbledhje dhe ushtrime}
Interpolimi supozon se të dhënat duhet të kalohen saktësisht; regresi pranon mbetje dhe kërkon një model të tyre. Zgjedhja nuk është thjesht numerike, por lidhet me kuptimin e të dhënave.
\begin{enumerate}
\item Vërtetoni formulën e mbetjes së Lagranzhit duke ndërtuar një funksion me $n+2$ zero.
\item Krahasoni nyjet e baraslarguara dhe të Çebishevit për funksionin Runge.
\item Nxirrni Richardson-in për një metodë të rendit të dytë dhe ndërtoni tabelën Romberg.
\item Përshtatni ligjin e tretë të Keplerit në shkallë lineare dhe logaritmike; diskutoni modelin e gabimit në secilin rast.
\item Formuloni regresin me shmangie minimale absolute si program linear.
\end{enumerate}
